\section{Binding a return obligation to its applicable equations}
\label{sec:requirement-successors}

A return obligation connects several different quantities: the number of goods
expected to be returned, the refund per unit, the carrying amount recoverable
from each return, and the balances already recorded. Checking their relationship
requires more than scalar bounds. The quantities must come from the intended
record, their units must combine correctly, and the equation must apply to the
state being examined. The successor requirement references make those three
steps explicit while preserving the first reference's supplied-input semantics.

Consider the retained conditional return example for legal entity P and its
2025 annual period. In its unsettled opening, expected returns are two units,
the refund is 20 dollars per unit, and the cost is 12 dollars per unit. Estimated
handling and deterioration reduce recovery by three and one dollars per unit,
respectively. These are supplied assumptions in a finite example. Its selected
opening equations are
\begin{align}
  L_{\mathrm{refund}} &= q p = 2\times20 = 40, \\
  A_{\mathrm{recovery}} &= q(c-h-d)
    = 2(12-3-1)=16.
\end{align}
Here $q$ is a quantity of goods; $p,c,h,d$ are dollars per unit; and both balances
are dollar amounts. The example's separate zero-tax-base conditions and supplied
tax rate of one quarter give the selected deferred-tax amounts 10 and 4. Those
premises belong to the retained conditional method. The equations alone establish
neither the premises nor the broader recognition and measurement treatment.

\subsection{Select the record before checking its amounts}

An input such as the refund balance must identify the exact record and field
from which it is read. Otherwise an equation could pass using a revised record,
another legal entity, or a consolidated amount that happens to share the same
number. The binder requires the caller to supply a plan, the selected method
definition, a selected-field schema, and the raw record bytes. It resolves
content identities from those supplied registries; it does not fetch files or
follow executable references.

The plan names the definition's canonical digest. Each named record has a raw
SHA256 digest and a schema digest, and each binding names an input, record and
JSON pointer. The input's source reference must equal
\code{sha256:RAW\_DIGEST\#POINTER}. The schema declares the field's complete type,
role, encoding and missing-path policy. Type, context and role must agree exactly
with the method input. All declared inputs must be mapped once, and every named
record must be used. Missing registry entries are unsupported; stale hashes,
duplicate mappings and undeclared field selections are invalid.

The raw decoder rejects duplicate keys, floating-point numerals, nonfinite
values, malformed UTF-8 and non-ASCII object keys. Unicode string values remain
available. Pointer escapes are checked before traversal, including when an
earlier part of the path is absent. A malformed pointer is an invalid request,
whereas an absent well-formed path follows the declared missing policy. Traversing
a scalar as though it were a collection is invalid. These distinctions prevent a
bad selector from being treated as merely unavailable financial information.

The three encoding choices also carry substantive meaning. A scalar is read as
a known value of its declared type. The explicitly selected legacy convention
treats the string \emph{unknown} as an unresolved value. A tagged binding supplies
its own known value or unknown reason. The missing-path policy is either invalid
or unresolved; unresolved paths receive the reason
\code{missing\_declared\_path}. None of these conventions supplies zero. A text
field using the scalar convention can contain the known word \emph{unknown}.

A successful binding returns the values and their record, pointer, schema,
type, role and source identity. The selected requirements are then assessed on
those values. Binding can therefore succeed while a method condition is invalid
or unresolved. Exact identity makes the calculation reproducible; the caller
still has to justify the schema's economic meaning and the source's authority.
The conformance examples that serialize supplied values into canonical records
identify those records as derived projections, rather than original filing bytes.

\subsection{Multiply quantities without erasing their units}

The dimensional successor, \code{MethodRequirements/2}, adds multiplication and
separates dimensions and scale from the context map. A rational type contains
its kind, context, dimension exponents and a positive rational scale. The
context identifies such matters as entity, legal book, component, period and
basis. A dimension name denotes a caller-declared quantity such as dollars or
return units. Zero exponents are omitted; exponents are bounded integers.

For two rational inputs in exactly the same context, let $a,b$ be their rational
coefficients, $s_a,s_b$ their scales and $d_{a,k},d_{b,k}$ the exponents of dimension
$k$. Multiplication produces
\[
 (ab,\ s_a s_b,\ \{k\mapsto d_{a,k}+d_{b,k}\}).
\]
Thus a return-unit quantity multiplied by dollars per return unit gives dollars:
the unit exponent cancels and the dollar exponent remains one. This rule defines
the local type calculation. It uses exact rational arithmetic and validates the
resulting type; it performs no inferred currency or scale conversion.

Addition, subtraction and comparison require exact agreement of context,
dimensions and scale. A dollar amount at scale one cannot be added directly to
a dollar amount at scale one thousand. Two different currencies do not match,
and a legal-book amount cannot be exchanged for a consolidated observation.
A conversion or allocation needs a separately specified relation. An older v1
context string that happens to contain a currency name is never promoted
automatically to this dimensional type.

The implementation lowers products to internal typed inputs, uses the preserved
scalar evaluator for the remaining operations, and restores the original
dependencies and unknown reasons in public findings. Internal product identities
and the reserved context keys used to encode dimensions are unavailable to caller
declarations in the dimensional versions. If a factor is unknown, the product
remains unknown even when another factor is zero. This follows the declared
conservative evaluation rule; it is not a symbolic proof about all admissible
histories.

\subsection{An inactive equation does not settle an unknown fact}

The unsettled return equation is inappropriate once the example's obligation
has been settled. \code{MethodRequirements/3} therefore adds a Boolean guard and
its exact dependency list to each requirement. The body predicate retains its
own dependency list. For the unsettled allocation, the guard is the negation of
the supplied settlement indicator; a separate settled requirement checks the
corresponding zero balance. The caller must author these conditions explicitly.

Guard evaluation has three distinct consequences. A true guard activates the
body's ordinary true, false or unknown result. A false guard marks the requirement
\code{not\_applicable}, even when the body is false or unknown. An unknown guard
leaves the requirement unresolved, even when its body happens to be true. The
result retains guard truth, body truth, activation and all original dependencies.
The interpreter checks structure and types in every branch before assessment;
an inactive branch cannot hide an unsupported operator or malformed value.

Some facts remain necessary independently of whether a particular equation is
active. A separate \code{required\_known} list declares those inputs. If the
refund balance is required but unknown, making the unsettled equation inactive
does not make the balance known. The result remains unresolved on that separate
condition. When all requirements are inactive and no other condition fails, the
summary is \code{no\_active\_requirements}. It does not assert that an accounting
method has been satisfied or can run.

For the zero-tax-base example, an equation's false guard also must not hide a
method limitation. Its zero-base scope checks are independently active: a known
nonzero base remains unsupported by the selected method, even though the
corresponding zero-base equation is inactive. This is why applicability, mandatory
knowledge and supported scope have separate declarations. Their adequacy remains
a question for the method's designer, supported by its source and state evidence.

\subsection{Keep versions explicit through authoring and exchange}

The new \code{MethodRequirementExchange/2} selects exactly one of the three
requirement profiles. It carries the full definition, selected digest, values
and available-capability names. The new text header names both the exchange and
requirement versions; input and constant declarations carry complete type
objects. The first two requirement versions use \code{require}; v3 uses
\code{guarded}, including separate guard and body dependencies and expressions.
Every v3 text document must contain exactly one \code{required\_known} list,
including an explicit empty list when intended. Omitting it is invalid.

The following fragment shows the guarded return relation after its quantity,
price, balance and settlement inputs have been declared. As in the preceding
reference, the display wraps one declaration for readability; the executable
text places each complete declaration on one line.
\begin{lstlisting}
guarded "refund_allocation" invalid unresolved
  ["settled"] ["expected","refund","refund_per_unit"]
  ["retained-unsettled-zero-tax-base-premises"]
  (not (ref "settled"))
  (eq (ref "refund")
      (mul (ref "expected") (ref "refund_per_unit")));
\end{lstlisting}
The complete document supplies the remaining declarations and its exact content
digest. Changing the guard, a type, a required-known list or a source reference
changes the selected definition. An unknown profile or incompatible declaration
is rejected. Profile names select a closed set of preserved evaluators; they
cannot name a caller-provided module or callback.

The corresponding \code{MethodInputBinding/2} requires a
\code{RequirementInputSchema/2} that explicitly names the expected requirement
profile. It retains the original full types and original plan, method, schema
and raw-record identities. Internally, a separately identified v1 representation
reuses the protected byte and pointer validation. This representation checks
binding only and is never assessed as the method. The original selected evaluator
receives the bound values, preserving dimensions, guards and unresolved reasons.

The recorded comparisons cover 32 cases, including 29 retained examples, and
59 combined unit-test methods. Full selected assessments agree across typed
definitions, text, exchange and binding; independently authored syntax and
adversarial changes supplement the roundtrips. These are deterministic checks
of the stated language fragment. Collection membership and reference order,
shared-book invariants, state effects, replay and reporting still need their
own complete method semantics. Selected requirements can explain why an input
is usable or why work must wait; they do not authorize the omitted transitions.
